\documentclass[a4paper,12pt]{article}

\usepackage[T2A]{fontenc}
\usepackage[utf8]{inputenc}
\usepackage[russian]{babel}
\usepackage{amsmath,amssymb,mathtools}
\usepackage{helvet}
\renewcommand{\familydefault}{\sfdefault}
\usepackage{icomma}
\usepackage[left=20mm,right=20mm,top=20mm,bottom=22mm]{geometry}
\usepackage{microtype}
\usepackage{tikz}
\usetikzlibrary{calc,arrows.meta,positioning,shapes.geometric}
\usepackage{tabularx,booktabs}
\usepackage{enumitem}
\usepackage{hyperref}
\usepackage{parskip}
\usepackage{fancyhdr}
\usepackage{setspace}
\usepackage{xcolor}

\definecolor{accent}{HTML}{008080}
\definecolor{darkaccent}{HTML}{004D4D}
\definecolor{textgray}{HTML}{333333}
\definecolor{softgray}{HTML}{F1F4F4}

\hypersetup{
    colorlinks=true,
    linkcolor=darkaccent,
    urlcolor=darkaccent
}

\color{textgray}
\onehalfspacing
\setlength{\parindent}{0pt}
\setlist[itemize]{leftmargin=7mm,itemsep=2pt,topsep=3pt}
\setlist[enumerate]{leftmargin=8mm,itemsep=5pt,topsep=4pt}
\renewcommand{\arraystretch}{1.35}

\pagestyle{fancy}
\fancyhf{}
\fancyhead[L]{\small Степени и показательные уравнения}
\fancyhead[R]{\small Екатерина Гнедкова}
\fancyfoot[C]{\small\thepage}
\renewcommand{\headrulewidth}{0.4pt}
\renewcommand{\headrule}{\hbox to\headwidth{\color{accent}\leaders\hrule height \headrulewidth\hfill}}

\newcommand{\sectionline}{%
    \par\noindent{\color{accent}\rule{\linewidth}{0.7pt}}\par
}

\newcommand{\important}[1]{%
    \par\smallskip
    {\color{darkaccent}\bfseries #1}
    \par\smallskip
}

\tikzset{
    flowstart/.style={
        ellipse,
        draw=darkaccent,
        line width=0.9pt,
        minimum width=35mm,
        minimum height=10mm,
        align=center,
        inner sep=4pt
    },
    flowproc/.style={
        rectangle,
        rounded corners=2mm,
        draw=accent,
        line width=0.9pt,
        text width=74mm,
        minimum height=12mm,
        align=center,
        inner sep=5pt
    },
    flowcond/.style={
        diamond,
        aspect=2.7,
        draw=darkaccent,
        line width=0.9pt,
        text width=45mm,
        minimum height=15mm,
        align=center,
        inner sep=2pt
    },
    flowarrow/.style={
        -{Stealth[length=3mm,width=2mm]},
        line width=0.9pt,
        color=darkaccent
    },
    flowlabel/.style={
        fill=white,
        inner sep=1.5pt,
        font=\small,
        text=darkaccent
    }
}

\begin{document}

%-----------------------------------------------------------------------
% Титульный лист
%-----------------------------------------------------------------------

\begin{titlepage}
\thispagestyle{empty}

\begin{center}
\vspace*{22mm}

{\color{darkaccent}\Large\bfseries МАТЕМАТИКА}

\vspace{17mm}

{\Huge\bfseries Чек-лист}

\vspace{7mm}

{\LARGE\color{accent}
Степенные преобразования\\[2mm]
и показательные уравнения}

\vspace{16mm}

{\large Подготовка к ЕГЭ по профильной математике}

\vfill

{\large Екатерина Гнедкова}

\vspace{9mm}

{\large \today}

\vspace{16mm}

{\small
Лёвин Артём Александрович\\
эксклюзивно для Екатерины Гнедковой}

\vspace{16mm}

\begin{tikzpicture}
\draw[accent,line width=1.2pt]
    (-6,0)
    .. controls (-3.8,0.8) and (-2.1,-0.7) .. (0,0)
    .. controls (2.2,0.7) and (4.0,-0.8) .. (6,0);
\fill[darkaccent] (0,0) circle (1.6pt);
\end{tikzpicture}

\end{center}
\end{titlepage}

%-----------------------------------------------------------------------
% Основной чек-лист
%-----------------------------------------------------------------------

\section*{\color{darkaccent}1. Цель преобразований}
\sectionline

Главная идея: представить все части выражения в степенном виде и получить
повторяющиеся основания, показатели или множители. После этого применяются
свойства степеней, вынесение общего множителя или сравнение показателей.

\important{Базовый маршрут}

\[
\boxed{\text{составные числа}}
\quad\longrightarrow\quad
\boxed{\text{корни}}
\quad\longrightarrow\quad
\boxed{\text{дроби}}
\quad\longrightarrow\quad
\boxed{\text{свойства степеней}}.
\]

Примеры перехода к степенному виду:
\[
64=2^6,\qquad
\sqrt[5]{a}=a^{1/5},\qquad
\frac{1}{b^4}=b^{-4}\quad (b\ne0).
\]

\section*{\color{darkaccent}2. Основные свойства степеней}
\sectionline

Для допустимых значений оснований и показателей:
\[
a^m\cdot a^n=a^{m+n},
\qquad
\frac{a^m}{a^n}=a^{m-n}\quad (a\ne0),
\]
\[
(a^m)^n=a^{mn},
\qquad
(ab)^n=a^n b^n,
\qquad
\left(\frac{a}{b}\right)^n=\frac{a^n}{b^n}\quad (b\ne0).
\]

Обратные преобразования также полезны:
\[
a^n b^n=(ab)^n,
\qquad
\frac{a^n}{b^n}=\left(\frac{a}{b}\right)^n.
\]

\important{Пример 1. Одинаковые основания}

\[
\frac{6^{\sqrt3}\cdot 7^{\sqrt3}}
     {42^{\sqrt3-1}}
=
\frac{42^{\sqrt3}}{42^{\sqrt3-1}}
=
42^{\sqrt3-(\sqrt3-1)}
=42.
\]

\important{Пример 2. Степень степени}

\[
\frac{(7m^6)^5+(11m^6)^5}
     {9m^{30}}
\]
в таком виде складывать слагаемые нельзя: коэффициенты также возводятся
в пятую степень. Если исходное выражение имеет вид
\[
\frac{7(m^6)^5+11(m^6)^5}{9m^{30}},
\]
то при \(m\ne0\):
\[
\frac{7m^{30}+11m^{30}}{9m^{30}}
=
\frac{18m^{30}}{9m^{30}}
=2.
\]

Запись скобок определяет, какая часть выражения возводится в степень.

\section*{\color{darkaccent}3. Деление и порядок действий}
\sectionline

Деление через двоеточие удобно сразу переписывать дробью:
\[
b^5:b^9\cdot b^6
=
\frac{b^5\cdot b^6}{b^9}
=
b^{5+6-9}
=
b^2,\qquad b\ne0.
\]

Здесь делителем является только \(b^9\). При записи
\[
b^5:(b^9\cdot b^6)
\]
делителем становится всё выражение в скобках:
\[
\frac{b^5}{b^{15}}=b^{-10}=\frac1{b^{10}},\qquad b\ne0.
\]

\important{Проверяйте}

\begin{itemize}
    \item границы числителя и знаменателя;
    \item наличие скобок;
    \item знак показателя;
    \item условие \(b\ne0\), если выражение содержит деление на \(b\).
\end{itemize}

\section*{\color{darkaccent}4. Показательные уравнения}
\sectionline

Показательным называется уравнение, в котором переменная входит в показатель
степени.

Если
\[
a^{f(x)}=a^{g(x)},\qquad a>0,\qquad a\ne1,
\]
то
\[
f(x)=g(x).
\]

\important{Пример 3. Приведение к одному основанию}

\[
2^{1-3x}=\frac1{64}.
\]
Поскольку
\[
64=2^6,\qquad \frac1{64}=2^{-6},
\]
получаем
\[
2^{1-3x}=2^{-6}
\quad\Longrightarrow\quad
1-3x=-6
\quad\Longrightarrow\quad
x=\frac73.
\]

\important{Пример 4. Составные числа и корень}

\[
25^{x+1}=\sqrt{125}.
\]
Переходим к основанию \(5\):
\[
(5^2)^{x+1}=(5^3)^{1/2},
\]
\[
5^{2x+2}=5^{3/2}.
\]
Следовательно,
\[
2x+2=\frac32,
\qquad
x=-\frac14.
\]

\section*{\color{darkaccent}5. Одинаковые показатели}
\sectionline

Иногда основания сразу сделать одинаковыми трудно. Тогда полезно собрать
степени с одинаковыми показателями:
\[
2^{x+3}=0{,}4\cdot 5^{x+3}.
\]
Переносим степенной множитель делением:
\[
\frac{2^{x+3}}{5^{x+3}}=0{,}4.
\]
Используем
\[
\frac{a^n}{b^n}=\left(\frac ab\right)^n,
\qquad
0{,}4=\frac25:
\]
\[
\left(\frac25\right)^{x+3}
=
\left(\frac25\right)^1.
\]
Так как \(\frac25>0\) и \(\frac25\ne1\), получаем
\[
x+3=1,
\qquad
x=-2.
\]

\section*{\color{darkaccent}6. Вынесение общей степени}
\sectionline

Формулы сложения показателей применяются к произведению. Для суммы или
разности используется вынесение общего множителя.

\important{Пример 5}

\[
6^{x+1}-6^x=180.
\]
Разделяем показатель:
\[
6^{x+1}=6^x\cdot6.
\]
Тогда
\[
6\cdot6^x-6^x=180,
\]
\[
6^x(6-1)=180,
\]
\[
6^x=36=6^2.
\]
Следовательно,
\[
x=2.
\]

\section*{\color{darkaccent}7. Как выделить неизвестную степень}
\sectionline

Если
\[
t^{3/5}=2^{1/5},
\]
требуется получить показатель \(1\) у переменной \(t\). Возводим обе части
в степень, обратную к \(\frac35\):
\[
\left(t^{3/5}\right)^{5/3}
=
\left(2^{1/5}\right)^{5/3}.
\]
Отсюда
\[
t=2^{1/3}=\sqrt[3]{2}.
\]

Для положительных величин общий приём имеет вид
\[
t^r=A
\quad\Longrightarrow\quad
t=A^{1/r},
\qquad r\ne0.
\]

При дробных показателях дополнительно проверяйте область допустимых значений.
Например, выражение \(x^{1/2}\) в действительных числах требует \(x\ge0\).

\section*{\color{darkaccent}8. Типичные ошибки}
\sectionline

\begin{itemize}
    \item Сложение показателей при сумме:
    \[
    a^m+a^n\ne a^{m+n}.
    \]
    \item Потеря скобок:
    \[
    (a^m)^n=a^{mn}.
    \]
    \item Неверное распределение степени:
    \[
    (a+b)^n\ne a^n+b^n
    \]
    в общем случае.
    \item Пропуск сокращения дроби:
    \[
    \frac{16}{20}=\frac45.
    \]
    \item Деление на выражение, которое может обращаться в нуль.
    \item Сравнение показателей при разных основаниях.
    \item Изменение числа или знака при переписывании условия.
    \item Нечёткая запись дробного показателя. Пишите
    \[
    a^{\frac15},
    \]
    чтобы показатель читался однозначно.
\end{itemize}

\clearpage

%-----------------------------------------------------------------------
% Блок-схема
%-----------------------------------------------------------------------



\clearpage

%-----------------------------------------------------------------------
% Тренировочный блок
%-----------------------------------------------------------------------

\section*{\color{darkaccent}Тренировочный блок}
\sectionline

Во всех заданиях выполняйте преобразования последовательно. Указывайте
ограничения, если выражение содержит переменную в знаменателе.

\subsection*{\color{darkaccent}Средний уровень}

\begin{enumerate}
    \item Вычислите:
    \[
    \frac{2^5\cdot2^{-2}}{2}.
    \]

    \item Упростите:
    \[
    \frac{(5^3)^2}{5^4}.
    \]

    \item Упростите выражение:
    \[
    \frac{7m^{12}+11m^{12}}{9m^{12}}.
    \]
\end{enumerate}

\subsection*{\color{darkaccent}Уровень выше среднего}

\begin{enumerate}[resume]
    \item Вычислите:
    \[
    \frac{2^{\sqrt5}\cdot4^{\sqrt5}}
         {8^{\sqrt5-1}}.
    \]

    \item Решите уравнение:
    \[
    2^{3x-1}=32.
    \]

    \item Решите уравнение:
    \[
    25^{x+1}=\sqrt{125}.
    \]

    \item Решите уравнение:
    \[
    \left(\frac37\right)^{2x-1}=\frac{49}{9}.
    \]

    \item Решите уравнение:
    \[
    \frac{2^{x+2}}{5^{x+2}}=\frac{16}{625}.
    \]

    \item Решите уравнение:
    \[
    6^{x+1}-6^x=180.
    \]
\end{enumerate}

\subsection*{\color{darkaccent}Сложный уровень}

\begin{enumerate}[resume]
    \item Найдите \(t\), если
    \[
    t^{3/5}=8.
    \]
\end{enumerate}

\clearpage

%-----------------------------------------------------------------------
% Ответы
%-----------------------------------------------------------------------

\section*{\color{darkaccent}Ответы}
\sectionline

\begin{table}[ht]
\centering
\caption{Ответы к тренировочному блоку}
\begin{tabularx}{\linewidth}{@{}cX@{}}
\toprule
\textbf{№} & \textbf{Ответ} \\
\midrule
1 & \(4\) \\

2 & \(25\) \\

3 & \(2\), при \(m\ne0\) \\

4 & \(8\) \\

5 & \(x=2\) \\

6 & \(x=-\dfrac14\) \\

7 & \(x=-\dfrac12\) \\

8 & \(x=2\) \\

9 & \(x=2\) \\

10 & \(t=32\) \\
\bottomrule
\end{tabularx}
\end{table}

\vfill

\begin{center}
{\color{darkaccent}\bfseries Итоговая памятка}

\vspace{3mm}

\begin{tabularx}{0.92\linewidth}{@{}>{\bfseries}p{39mm}X@{}}
\toprule
Сигнал в задаче & Полезное действие \\
\midrule
Составное число & Разложите его на простые множители или представьте степенью. \\

Корень & Запишите корень как дробную степень. \\

Дробь & Сократите её и при необходимости используйте отрицательный показатель. \\

Одинаковые основания & Сравните показатели. \\

Одинаковые показатели & Объедините основания. \\

Повторяющийся фрагмент & Вынесите общий множитель. \\

Степень у неизвестной & Используйте обратный показатель. \\
\bottomrule
\end{tabularx}
\end{center}

\end{document}